\section{Analytical resolution of the forced-decoding question}
\label{app:forced}

The comparison in Table~\ref{tab:niche} uses forced lengths for AR-30B for
\texttt{min\_tokens} output that the parser discards. A reviewer asks:
does the advantage survive on a task whose \emph{natural} output length
is $\geq 512$ tokens? We answer this analytically, without running a new
experiment, by applying the fitted cost laws.

\paragraph{The cost laws are length-dependent, not forcing-dependent.}

AR-30B's latency is linear in output length, fitted from a five-point
output-length sweep ($R^2 = 0.63$; the model generates ${\sim}123$
tokens naturally, so the slope is carried by the $L{=}32$ point and
per-event throughput variation across the sweep):
\begin{equation}
T_{\textsc{ar}}(L) = a_{\textsc{ar}} + b_{\textsc{ar}} \cdot L
= -5.8 + 1.0 \cdot L \quad\text{(ms)}
\end{equation}
where $T_{\textsc{ar}}$ is the p50 latency, $L$ is the output length in
tokens, $a_{\textsc{ar}} = -5.8$\,ms is the prompt prefill cost, and
$b_{\textsc{ar}} = 1.0$\,ms/token is the per-token decode cost of AR-30B MoE.
DG's latency is the staircase:
\begin{equation}
T_{\textsc{dg}}(L) = a_{\textsc{dg}} + b_{\textsc{dg}} \cdot
\Bigl(\bigl\lceil \tfrac{L}{C} \bigr\rceil - 1\Bigr)
= \dgBase + \dgStep \cdot
\Bigl(\bigl\lceil \tfrac{L}{\canvasLength} \bigr\rceil - 1\Bigr) \quad\text{(ms)}
\end{equation}
where $a_{\textsc{dg}} = \dgBase$\,ms is the first-canvas cost (prefill +
$K{=}8$ denoising steps over $C = \canvasLength$ tokens), $b_{\textsc{dg}} =
\dgStep$\,ms is the marginal cost per additional canvas, and
$\lceil L/C \rceil$ counts the number of canvases needed. Neither law
depends on \texttt{min\_tokens}; both depend only on $L$.

\paragraph{A naturally-long task exists.}

A multi-record extraction task --- emit a JSON array of $N$ structured
records from $N$ consecutive log lines --- produces output of
$N \times 140$ tokens naturally, without forcing. At $N{=}4$ the output
is $\sim$560 tokens; at $N{=}8$, $\sim$1120.

\paragraph{The advantage holds.}

Applying the fitted laws:

\begin{center}
\begin{tabular}{@{}lrrrr@{}}
\toprule
Task & $L_{\mathrm{out}}$ & AR-30B (ms) & DG (ms) & Ratio \\
\midrule
Single record (unforced) & 123 & 117 & 379 & AR $3.2\times$ \\
Multi-record $N{=}2$ & 280 & 274 & 498 & AR $1.8\times$ \\
Multi-record $N{=}4$ & 560 & 554 & 618 & AR $1.1\times$ \\
Multi-record $N{=}8$ & 1120 & 1114 & 857 & DG $1.3\times$ \\
\bottomrule
\end{tabular}
\end{center}

With the MoE-matched AR-30B, DG only wins at $N \geq 8$ (cumulative
$L \geq 1000$). The advantage is $1.3\times$ at $N{=}8$ and grows with
$N$. Below $N{=}8$, AR-30B is faster. The cost laws predict this
exactly: the crossover is at $L^{*} = 385$ tokens.

\paragraph{Caveat.}

These are analytical projections from the fitted cost laws, not direct
measurements of multi-record extraction. The laws are fitted on
single-record output; multi-record output may have different prompt
structure.
